\section*{Chapitre \ref{ch:g11:functional-groups} --- Groupes caractéristiques et familles}

\begin{solution}{exo:g11:functional-groups:1}
Hydroxyle, \omterm{def:g11:functional-groups:alcohol}{alcool}~; carbonyle à l'intérieur de la chaîne, \omterm{def:g11:functional-groups:carbonyl}{cétone}~;
carboxyle, \omterm{def:g11:functional-groups:carboxylic-acid}{acide carboxylique}~; groupe \omterm{def:g11:functional-groups:carboxylic-acid}{ester}, \omterm{def:g11:functional-groups:carboxylic-acid}{ester}~; groupe amino
\ce{-NH2}, \omterm{def:g11:functional-groups:amine}{amine}.
\end{solution}

\begin{solution}{exo:g11:functional-groups:2}
Méthanol~; propan-1-ol~; acide méthanoïque~; propanal~; propanone.
\end{solution}

\begin{solution}{exo:g11:functional-groups:3}
\omterm{def:g11:functional-groups:alcohol}{Alcool}~; \omterm{def:g11:functional-groups:carbonyl}{aldéhyde}~; \omterm{def:g11:functional-groups:carboxylic-acid}{ester}~; \omterm{def:g11:functional-groups:amine}{amide}~; \omterm{def:g11:functional-groups:halogenoalkane}{halogénoalcane}~; \omterm{def:g11:functional-groups:halogenoalkane}{alcène}.
\end{solution}

\begin{solution}{exo:g11:functional-groups:4}
Propanal \ce{CH3-CH2-CHO}, propanone \ce{CH3-CO-CH3}. Le propanal est
l'\omterm{def:g11:functional-groups:carbonyl}{aldéhyde}~: le carbone de son \ce{C=O} est en bout de chaîne, lié à un
\omterm{def:g7:atoms-and-molecules:atom}{atome} d'hydrogène~; dans la propanone, il est entre deux \omterm{def:g7:atoms-and-molecules:atom}{atomes} de
carbone.
\end{solution}

\begin{solution}{exo:g11:functional-groups:5}
Primaire (le carbone qui porte le \ce{-OH} a un carbone voisin)~;
secondaire (deux)~; tertiaire (trois).
\end{solution}

\begin{solution}{exo:g11:functional-groups:6}
\setchemfig{atom sep=1.5em}
Pentan-3-ol \chemfig{-[:30]-[:-30](-[:-90]OH)-[:30]-[:-30]}~;
acide 2-méthylbutanoïque \chemfig{-[:30]-[:-30](-[:-90])-[:30](=[:90]O)-[:-30]OH}~;
hexan-2-one \chemfig{-[:30](=[:90]O)-[:-30]-[:30]-[:-30]-[:30]}~;
éthanoate de propyle \chemfig{-[:30](=[:90]O)-[:-30]O-[:30]-[:-30]-[:30]}.
\end{solution}

\begin{solution}{exo:g11:functional-groups:7}
\ce{C3H6O}. Oui, ce sont des \omterm{def:g11:organic-skeletons:isomer}{isomères}, mais avec des groupes
différents~: un \omterm{def:g11:functional-groups:carbonyl}{aldéhyde} et une \omterm{def:g11:functional-groups:carbonyl}{cétone} (deux groupes carbonyle, à des
places différentes).
\end{solution}

\begin{solution}{exo:g11:functional-groups:8}
\ce{HCOO-CH2-CH3}, méthanoate d'éthyle~; \ce{CH3-COO-CH2-CH2-CH3},
éthanoate de propyle.
\end{solution}

\begin{solution}{exo:g11:functional-groups:9}
Éthanoate de méthyle (acide éthanoïque, méthanol)~; propanoate d'éthyle
(acide propanoïque, éthanol)~; méthanoate de méthyle (acide méthanoïque,
méthanol).
\end{solution}

\begin{solution}{exo:g11:functional-groups:10}
\setchemfig{atom sep=1.5em}
\begin{center}
\small
\begin{tabular}{cccc}
\chemfig{-[:30]-[:-30]-[:30]-[:-30]OH} &
\chemfig{-[:30](-[:90]OH)-[:-30]-[:30]} &
\chemfig{-[:30](-[:90])-[:-30]-[:30]OH} &
\chemfig{-[:0](-[:90]OH)(-[:-90])-[:0]} \\
butan-1-ol & butan-2-ol & 2-méthylpropan-1-ol & 2-méthylpropan-2-ol \\
primaire & secondaire & primaire & tertiaire
\end{tabular}
\end{center}
\end{solution}

\begin{solution}{exo:g11:functional-groups:11}
Un \omterm{def:g11:functional-groups:amine}{amide} (l'éthanamide). \emph{-ène}. Les \omterm{def:g11:functional-groups:carboxylic-acid}{acides carboxyliques}
(\ce{C=O} et \ce{-OH} sur le même carbone), les \omterm{def:g11:functional-groups:carboxylic-acid}{esters} (\ce{C=O} et
\ce{-O-}) et les \omterm{def:g11:functional-groups:amine}{amides} (\ce{C=O} et \ce{N})~:
trois familles.
\end{solution}

\begin{solution}{exo:g11:functional-groups:12}
L'aspirine~: un groupe carboxyle (\omterm{def:g11:functional-groups:carboxylic-acid}{acide carboxylique}) et un groupe \omterm{def:g11:functional-groups:carboxylic-acid}{ester}.
Le paracétamol~: un \ce{-OH} de phénol sur le cycle et un groupe \omterm{def:g11:functional-groups:amine}{amide}
\ce{NH-CO}.
\end{solution}

\begin{solution}{exo:g11:functional-groups:13}
\ce{C2H6O}. Seul l'éthanol, avec son \ce{O-H}, forme des \omterm{def:g11:polarity-and-cohesion:hydrogen-bond}{liaisons
hydrogène} entre ses \omterm{def:g7:atoms-and-molecules:molecule}{molécules} (le méthoxyméthane n'a pas d'hydrogène sur
son oxygène). L'éthanol bout le plus haut, à \qty{78}{\celsius}~; le
méthoxyméthane à \qty{-25}{\celsius}.
\end{solution}

\begin{solution}{exo:g11:functional-groups:14}
Un groupe hydroxyle et un groupe carboxyle. La chaîne compte trois
carbones, le carbone de l'acide étant le carbone 1 et le \ce{-OH} étant
sur le carbone 2~: acide 2-hydroxypropanoïque. La glycine~: acide
2-aminoéthanoïque (souvent écrit acide aminoéthanoïque).
\end{solution}

\begin{solution}{exo:g11:functional-groups:15}
L'éthanol (\ce{CH3-CH2-OH}, \ce{C2H6O}) est un \omterm{def:g11:functional-groups:alcohol}{alcool}~; l'éthanal
(\ce{CH3-CHO},
\ce{C2H4O}), un \omterm{def:g11:functional-groups:carbonyl}{aldéhyde}~; l'acide éthanoïque (\ce{CH3-COOH}, \ce{C2H4O2}),
un \omterm{def:g11:functional-groups:carboxylic-acid}{acide carboxylique}. De l'éthanol à l'éthanal, deux \omterm{def:g7:atoms-and-molecules:atom}{atomes} d'hydrogène
sont perdus~; de l'éthanal à l'acide éthanoïque, un \omterm{def:g7:atoms-and-molecules:atom}{atome} d'oxygène est
gagné.
\end{solution}

\begin{solution}{pb:g11:functional-groups:1}
\textbf{1.} A~: groupe \omterm{def:g11:functional-groups:carboxylic-acid}{ester}~; B~: hydroxyle~; C~: carboxyle~; D~:
carbonyle en bout de chaîne~; E~: groupe \omterm{def:g11:functional-groups:carboxylic-acid}{ester}.

\textbf{2.} A \omterm{def:g11:functional-groups:carboxylic-acid}{ester}, B \omterm{def:g11:functional-groups:alcohol}{alcool}, C \omterm{def:g11:functional-groups:carboxylic-acid}{acide carboxylique}, D \omterm{def:g11:functional-groups:carbonyl}{aldéhyde}, E \omterm{def:g11:functional-groups:carboxylic-acid}{ester}.

\textbf{3.} A (banane) et E (ananas)~: des odeurs fruitées.

\textbf{4.} Primaire~: le carbone qui porte le \ce{-OH} est lié à un seul
carbone.

\textbf{5.} \ce{C6H12O}~; par exemple l'hexan-2-one,
\ce{CH3-CO-CH2-CH2-CH2-CH3}.

\textbf{6.} La chaîne qui contient le \ce{C-OH} compte quatre carbones,
numérotés à partir du \ce{-OH}, avec un méthyle sur le carbone 3~:
3-méthylbutan-1-ol.

\textbf{7.} C~: acide éthanoïque~; D~: hexanal.

\textbf{8.} Butanoate d'éthyle, apparenté à l'acide butanoïque et à
l'éthanol.

\textbf{9.} \emph{Éthanoate}~: la partie acide \ce{CH3-CO}, de l'acide
éthanoïque~; \emph{3-méthylbutyle}~: la partie \omterm{def:g11:functional-groups:alcohol}{alcool}, du
3-méthylbutan-1-ol.

\textbf{10.} C (acide éthanoïque) et B (3-méthylbutan-1-ol).

\textbf{11.} A \ce{C7H14O2}~; B \ce{C5H12O}~; C \ce{C2H4O2}.

\textbf{12.} $\ce{C2H4O2} + \ce{C5H12O}$ contient \ce{C7H16O3}, soit un
\ce{H2O} de plus que \ce{C7H14O2}~: de l'eau.

\textbf{13.} \ce{C2H4O2 + C5H12O -> C7H14O2 + H2O}~: C 7 = 7, H 16 = 16,
O 3 = 3.

\textbf{14.} $M(\text{B}) = 5 \times 12.0 + 12 \times 1.0 + 16.0 =
\qty{88.0}{g/mol}$~; $M(\text{C}) = \qty{60.0}{g/mol}$.

\textbf{15.} $60.0 + 88.0 - 18.0 = \qty{130.0}{g/mol}$.

\textbf{16.} $7 \times 12.0 + 14 \times 1.0 + 2 \times 16.0 =
\qty{130.0}{g/mol}$~: la même valeur.

\textbf{17.} E, \ce{C6H12O2}~: $6 \times 12.0 + 12 \times 1.0 + 2 \times
16.0 = \qty{116.0}{g/mol}$.

\textbf{18.} \ce{C7H14O2}~: la même formule que A, donc un \omterm{def:g11:organic-skeletons:isomer}{isomère}, de
même \omterm{def:g10:the-mole:molar-mass}{masse molaire}. Une \omterm{def:g10:the-mole:molar-mass}{masse molaire} seule ne permet pas de distinguer
des \omterm{def:g11:organic-skeletons:isomer}{isomères}.

\textbf{19.} A et l'acide heptanoïque ont les mêmes \omterm{def:g7:atoms-and-molecules:atom}{atomes} mais des
groupes différents, un \omterm{def:g11:functional-groups:carboxylic-acid}{ester} et un \omterm{def:g11:functional-groups:carboxylic-acid}{acide carboxylique}~: l'un sent la
banane, l'autre le rance. C'est le groupe, et sa place, qui décide.

\textbf{20.} \qty{130.0}{g/mol}.
\end{solution}
